\documentclass[../cours_tle.tex]{subfiles}
\begin{document}
\setcounter{chapter}{13}
\chapter{Dériver, composer et contrôler les variations}
\label{ch:14}
\begin{objectif}
Spécialité. Retrouver les règles de dérivation et comprendre leur portée. Prérequis : dérivation de première et chapitre 13.
\end{objectif}
\section{La dérivée est une approximation locale}
$f$ est dérivable en $a$ si $(f(a+h)-f(a))/h$ tend vers un réel $f'(a)$.
Cela équivaut à
$f(a+h)=f(a)+hf'(a)+h\varepsilon(h)$ avec $\varepsilon(h)\to0$.
La droite $y=f(a)+f'(a)(x-a)$ est la tangente.
Cette écriture prouve que la dérivabilité entraîne la continuité. La réciproque est fausse : $|x|$ est continue en $0$, mais ses pentes à gauche et à droite sont $-1$ et $1$.
\section{Règles de calcul, avec leurs justifications}
La linéarité découle de celle des limites. Pour le produit, écrire
\[
\frac{u(a+h)v(a+h)-u(a)v(a)}h
=u(a+h)\frac{v(a+h)-v(a)}h+v(a)\frac{u(a+h)-u(a)}h.
\]
En passant à la limite, $(uv)'=u'v+uv'$. Pour $v\ne0$, le taux d'accroissement de $1/v$ donne $-(v')/v^2$ ; on en déduit $(u/v)'=(u'v-uv')/v^2$.
Par récurrence sur la règle du produit, $(x^n)'=nx^{n-1}$ pour $n\ge1$, puis pour les entiers négatifs avec l'inverse, sur $x\ne0$.
Pour $x>0$, le conjugué donne $(\sqrt x)'=1/(2\sqrt x)$.
\section{Dérivée d'une composition}
\begin{theoreme}
Si $u$ est dérivable en $a$ et $v$ dérivable en $u(a)$, alors
$(v\circ u)'(a)=v'(u(a))u'(a)$.
\end{theoreme}
\begin{demonstration}
Écrire $u(a+h)-u(a)=h(u'(a)+\eta(h))$, avec $\eta(h)\to0$.
Pour $k$ proche de $0$, $v(u(a)+k)-v(u(a))=k(v'(u(a))+\varepsilon(k))$ ; poser $\varepsilon(0)=0$.
En remplaçant $k$ par $u(a+h)-u(a)$ et divisant par $h$, on obtient la limite annoncée, même si $k=0$ pour certains $h$.
\end{demonstration}
Ainsi $(e^{u})'=u'e^u$, $(u^n)'=nu'u^{n-1}$ et, si $u>0$, $(\sqrt u)'=u'/(2\sqrt u)$.
Par exemple $(e^{-x^2})'=-2xe^{-x^2}$ : oublier la dérivée intérieure changerait les variations.
\section{Pourquoi le signe de la dérivée commande les variations}
\approfondissement
Si $f$ est continue sur $[a,b]$ et dérivable sur $]a,b[$, il existe $c\in]a,b[$ tel que
$f(b)-f(a)=f'(c)(b-a)$ : théorème des accroissements finis.
\begin{demonstration}
Soustraire à $f$ sa droite sécante donne $g(a)=g(b)$. Une fonction continue atteint ses bornes sur un segment (preuve au chapitre 36). Si $g$ n'est pas constante, un maximum ou minimum est intérieur. À un extremum intérieur dérivable, les taux d'accroissement de signes opposés imposent $g'=0$ ; donc $f'(c)$ est la pente de la sécante. Si $g$ est constante, tout point convient.
\end{demonstration}
Si $f'\ge0$ sur un intervalle, appliquer ce théorème entre deux points donne $f(b)\ge f(a)$ ; si $f'>0$, l'inégalité est stricte.
Réciproquement, les quotients d'une fonction croissante dérivable sont positifs ou nuls, donc sa dérivée l'est aussi.
Une dérivée nulle sur un intervalle donne une constante. L'hypothèse d'intervalle est nécessaire : sur $\R\setminus\{0\}$, on peut avoir deux constantes différentes.
\section{Optimiser et étudier une fonction}
Pour un rectangle de périmètre $20$, les côtés sont $x$ et $10-x$, $0<x<10$. Son aire $A(x)=10x-x^2$ a dérivée $10-2x$ ; le maximum est $25$, atteint pour le carré $x=5$.
Un point où $f'=0$ n'est pas forcément un extremum : $x^3$ en $0$.
Pour une étude complète : domaine, limites, dérivée, signe, variations, valeurs remarquables, puis graphique cohérent. Sur un segment, comparer aussi les valeurs aux bornes.
\section{Complément : dériver plusieurs fois}
On note $f^{(0)}=f$, $f^{(1)}=f'$ et $f^{(n+1)}=(f^{(n)})'$ lorsque ces fonctions existent. Pour $f(x)=e^{ax}$, une récurrence utilisant la chaîne donne $f^{(n)}(x)=a^ne^{ax}$. Pour $x^m$, les dérivations successives donnent $m(m-1)\cdots(m-n+1)x^{m-n}$ si $n\le m$, puis zéro si $n>m$. Pour le sinus, le cycle des quatre dérivées donne $\sin^{(n)}x=\sin(x+n\pi/2)$.
Ces calculs prépareront les développements de Taylor ; aucune approximation par un polynôme d'ordre élevé ne sera utilisée sans un contrôle de son reste.
\begin{finchapitre}
\section{Utilité et suite du parcours}
La formule de composition deviendra la règle de chaîne pour des applications de plusieurs variables. Les accroissements finis donnent des bornes d'erreur et la condition de contraction.
\begin{retenir}
Une dérivée exprime une propriété locale. Pour conclure sur un intervalle entier, contrôler son signe partout et vérifier les bornes.
\end{retenir}
\end{finchapitre}
\end{document}
